
    %%%%%%%%%%%%%%%%%%%%%%%%%%%%%%%%%%%%%%%%%%%%%%%%%%%%%%%%%%%%%%%%%%%%%%%%%%%%%%%%
    %%% Classe do documento
    \documentclass[12pt, addpoints]{exam}
    \UseRawInputEncoding

    %%%%%%%%%%%%%%%%%%%%%%%%%%%%%%%%%%%%%%%%%%%%%%%%%%%%%%%%%%%%%%%%%%%%%%%%%%%%%%%%
    % preambulo.tex
    %
    % Arquivo de configuração de packages para uso o LaTeX, conforme minhas
    % preferências e modelos pessoais.
    %
    % Para maiores informações, visite:
    %    
    %
    % NÃO ALTERE SE NÃO SOUBER O QUE ESTÁ FAZENDO!
    %%%%%%%%%%%%%%%%%%%%%%%%%%%%%%%%%%%%%%%%%%%%%%%%%%%%%%%%%%%%%%%%%%%%%%%%%%%%%%%%


    %%%%%%%%%%%%%%%%%%%%%%%%%%%%%%%%%%%%%%%%%%%%%%%%%%%%%%%%%%%%%%%%%%%%%%%%%%%%%%%%
    %%% Estruturas de controle:
    \input{static/utils/controles}

    %%%%%%%%%%%%%%%%%%%%%%%%%%%%%%%%%%%%%%%%%%%%%%%%%%%%%%%%%%%%%%%%%%%%%%%%%%%%%%%%
    %%% Compilação condicional em PDF:
    \input{static/utils/pdf}

    %%%%%%%%%%%%%%%%%%%%%%%%%%%%%%%%%%%%%%%%%%%%%%%%%%%%%%%%%%%%%%%%%%%%%%%%%%%%%%%%
    %%% Configurações de layout da página:
    \input{static/utils/layout}

    %%%%%%%%%%%%%%%%%%%%%%%%%%%%%%%%%%%%%%%%%%%%%%%%%%%%%%%%%%%%%%%%%%%%%%%%%%%%%%%%
    %%% Configurações para autores:
    \input{static/utils/autores}

    %%%%%%%%%%%%%%%%%%%%%%%%%%%%%%%%%%%%%%%%%%%%%%%%%%%%%%%%%%%%%%%%%%%%%%%%%%%%%%%%
    %%% Cabeçalho e rodapé:
    %%% Atenção! É MUITO DIFÍCIL ajustar apropriadamente os running headers
    %%% e running footers da primeira vez. Você pode confiar no LaTeX e deixar que
    %%% ele faça o ajuste automaticamente ou pode quebrar a cabeça e
    %%% tentar melhorar algo que o LaTeX já faz muito bem, mesmo que não fique
    %%% exatamente do jeito que você quer. O que você prefere? Se quiser ajustar
    %%% manualmente, descomente a linha abaixo e faça as configurações no arquivo
    %%% "utils/headings.tex".
    %\input{static/utils/headings}

    %%%%%%%%%%%%%%%%%%%%%%%%%%%%%%%%%%%%%%%%%%%%%%%%%%%%%%%%%%%%%%%%%%%%%%%%%%%%%%%%
    %%% Configuração para as fontes do documento:
    %%% Se você quiser usar fontes próprias ou do sistema, precisará ajustar as
    %%% configurações no arquivo "utils/fontes.tex".
    %%% ATENÇÃO: por padrão eu utilizo XeLaTeX com fontes proprietárias projetadas
    %%% por Matthew Butterick, adquiridas em https://mbtype.com, que não podem ser
    %%% distribuídas devido à licença de utilização. Se você usar XeLaTeX, você
    %%% DEVERÁ OBRIGATORIAMENTE ajustar as fontes no arquivo "utils/fontes.tex".
    \input{static/utils/fontes}

    %%%%%%%%%%%%%%%%%%%%%%%%%%%%%%%%%%%%%%%%%%%%%%%%%%%%%%%%%%%%%%%%%%%%%%%%%%%%%%%%
    %%% Sumário:
    \input{static/utils/toc}

    %%%%%%%%%%%%%%%%%%%%%%%%%%%%%%%%%%%%%%%%%%%%%%%%%%%%%%%%%%%%%%%%%%%%%%%%%%%%%%%%
    %%% Matemática:
    \input{static/utils/matematica}

    %%%%%%%%%%%%%%%%%%%%%%%%%%%%%%%%%%%%%%%%%%%%%%%%%%%%%%%%%%%%%%%%%%%%%%%%%%%%%%%%
    %%% Notas de rodapé:
    \input{static/utils/footnotes}

    %%%%%%%%%%%%%%%%%%%%%%%%%%%%%%%%%%%%%%%%%%%%%%%%%%%%%%%%%%%%%%%%%%%%%%%%%%%%%%%%
    %%% Referências cruzadas, links, citações:
    \input{static/utils/crossrefs}

    %%%%%%%%%%%%%%%%%%%%%%%%%%%%%%%%%%%%%%%%%%%%%%%%%%%%%%%%%%%%%%%%%%%%%%%%%%%%%%%%
    %%% Configurações para os hiperlinks em PDF:
    \input{static/utils/pdfconfig}

    %%%%%%%%%%%%%%%%%%%%%%%%%%%%%%%%%%%%%%%%%%%%%%%%%%%%%%%%%%%%%%%%%%%%%%%%%%%%%%%%
    %%% Glossário e índice remissivo:
    \input{static/utils/glossindex}

    %%%%%%%%%%%%%%%%%%%%%%%%%%%%%%%%%%%%%%%%%%%%%%%%%%%%%%%%%%%%%%%%%%%%%%%%%%%%%%%%
    %%% Referências bibliográficas:
    \input{static/utils/bibliografia}

    %%%%%%%%%%%%%%%%%%%%%%%%%%%%%%%%%%%%%%%%%%%%%%%%%%%%%%%%%%%%%%%%%%%%%%%%%%%%%%%%
    %%% Cores:
    \input{static/utils/cores}

    %%%%%%%%%%%%%%%%%%%%%%%%%%%%%%%%%%%%%%%%%%%%%%%%%%%%%%%%%%%%%%%%%%%%%%%%%%%%%%%%
    %%% Computação:
    \input{static/utils/computacao}

    %%%%%%%%%%%%%%%%%%%%%%%%%%%%%%%%%%%%%%%%%%%%%%%%%%%%%%%%%%%%%%%%%%%%%%%%%%%%%%%%
    %%% Gráficos:
    \input{static/utils/graficos}

    %%%%%%%%%%%%%%%%%%%%%%%%%%%%%%%%%%%%%%%%%%%%%%%%%%%%%%%%%%%%%%%%%%%%%%%%%%%%%%%%
    %%% Ícones:
    \input{static/utils/icones}

    %%%%%%%%%%%%%%%%%%%%%%%%%%%%%%%%%%%%%%%%%%%%%%%%%%%%%%%%%%%%%%%%%%%%%%%%%%%%%%%%
    %%% Blocos:
    \input{static/utils/blocos}

    %%%%%%%%%%%%%%%%%%%%%%%%%%%%%%%%%%%%%%%%%%%%%%%%%%%%%%%%%%%%%%%%%%%%%%%%%%%%%%%%
    %%% Boxes coloridos:
    \input{static/utils/colorbox}

    %%%%%%%%%%%%%%%%%%%%%%%%%%%%%%%%%%%%%%%%%%%%%%%%%%%%%%%%%%%%%%%%%%%%%%%%%%%%%%%%
    %%% Tabelas:
    \input{static/utils/tabelas}

    %%%%%%%%%%%%%%%%%%%%%%%%%%%%%%%%%%%%%%%%%%%%%%%%%%%%%%%%%%%%%%%%%%%%%%%%%%%%%%%%
    %%% Ambientes de listas:
    \input{static/utils/listas}

    %%%%%%%%%%%%%%%%%%%%%%%%%%%%%%%%%%%%%%%%%%%%%%%%%%%%%%%%%%%%%%%%%%%%%%%%%%%%%%%%
    %%% Ambientes floats e similares:
    \input{static/utils/floats}

    %%%%%%%%%%%%%%%%%%%%%%%%%%%%%%%%%%%%%%%%%%%%%%%%%%%%%%%%%%%%%%%%%%%%%%%%%%%%%%%%
    %%% Ajustes para a classe EXAM:
    \input{static/utils/exam}

    %%%%%%%%%%%%%%%%%%%%%%%%%%%%%%%%%%%%%%%%%%%%%%%%%%%%%%%%%%%%%%%%%%%%%%%%%%%%%%%%
    %%% Ajustes para textos:
    \input{static/utils/texto}

    %%%%%%%%%%%%%%%%%%%%%%%%%%%%%%%%%%%%%%%%%%%%%%%%%%%%%%%%%%%%%%%%%%%%%%%%%%%%%%%%
    %%% Ambientes, aliases, comandos e customizações em geral para serem
    %%% utilizadas nos documentos. Defina aqui qualquer customização específica
    %%% para seus documentos!
    \input{static/utils/customizacoes}

    %%%%%%%%%%%%%%%%%%%%%%%%%%%%%%%%%%%%%%%%%%%%%%%%%%%%%%%%%%%%%%%%%%%%%%%%%%%%%%%%
    %%% Ajuste do layout, espaçamento de linhas e etc.:
    \geometry{a4paper, portrait, left=2.0cm, right=2.0cm, top=2.0cm, bottom=2.0cm}
    \singlespacing

    %%%%%%%%%%%%%%%%%%%%%%%%%%%%%%%%%%%%%%%%%%%%%%%%%%%%%%%%%%%%%%%%%%%%%%%%%%%%%%%%
    %%% Configurações de cabeçalho e rodapé (não use fancyhdr pois ocorre conflito):
    \pagestyle{headandfoot}
    \runningheadrule
    \firstpageheader{}{}{}
    \runningheader{Sem disciplina}{}{Avaliação}
    \firstpagefooter{}{}{Página \thepage\ de \numpages}
    \runningfooter{}{}{Página \thepage\ de \numpages}
    \runningfootrule

    %%%%%%%%%%%%%%%%%%%%%%%%%%%%%%%%%%%%%%%%%%%%%%%%%%%%%%%%%%%%%%%%%%%%%%%%%%%%%%%%
    %%% Configurações para as propriedades do PDF:
    \hypersetup{
    hidelinks,           % Comente para web, descomente para impressão
    colorlinks=true,      % True para web, False para impressão
    pdftitle=Prova Par turma noite: Avaliação,
    pdfauthor=Professor,
    pdfsubject=Sem disciplina,
    pdfkeywords=['Sem tag'],
    pdfinfo={
        CreationDate={}, % Ex.: D:AAAAMMDDHH24MISS
        ModDate={}       % Ex.: D:AAAAMMDDHH24MISS
    }
    }
    \input{static/utils/pdfconfig}

    %%%%%%%%%%%%%%%%%%%%%%%%%%%%%%%%%%%%%%%%%%%%%%%%%%%%%%%%%%%%%%%%%%%%%%%%%%%%%%%%
    %%% Começa o documento
    \begin{document}

    %%%%%%%%%%%%%%%%%%%%%%%%%%%%%%%%%%%%%%%%%%%%%%%%%%%%%%%%%%%%%%%%%%%%%%%%%%%%%%%%
    %%% Folha de instruções padronizadas da para a prova
    \pdfbookmark[1]{Instruções}{instrucoes}
    \begin{figure}[H]
        \begin{center}
            \includegraphics[scale=0.3]{{static/imgs/logo-.png}}
        \end{center}	
    \end{figure}

    \vspace{-1cm}

    \begin{table}[H]
        \centering
        \begin{tabular}{|llll|}
            \hline
            \multicolumn{2}{|l|}{\textbf{Disciplina:} Sem disciplina} & 
            \multicolumn{1}{l|}{\multirow{3}{*}{\textbf{\makecell{Nota:\\ \,\\ \,}}}} & 
            \multirow{3}{*}{\textbf{\makecell{Coordenador:\\ \,\\ \,}}} \\ 
            \cline{1-2}
            \multicolumn{2}{|l|}{\textbf{Professor:} Professor \hspace{4.72cm} } & 
            \multicolumn{1}{l|}{} & \\ 
            \cline{1-2}
            \multicolumn{2}{|l|}{\textbf{Aluno:}} & 
            \multicolumn{1}{l|}{} & \\ 
            \hline
            \multicolumn{1}{|l|}{\textbf{Turma:} \hspace{6.3cm} } & 
            \multicolumn{1}{l|}{\textbf{Semestre:}} & 
            \multicolumn{2}{l|}{\textbf{Valor:} 10 pontos} \\ 
            \hline
            \multicolumn{1}{|l|}{\textbf{Data:}} & 
            \multicolumn{3}{l|}{\textbf{Avaliação:}} \\ 
            \hline
        \end{tabular}
    \end{table}

    \begin{center}
        \fbox{\setlength\fboxsep{0.3cm} \fbox{\parbox{15.4cm}{
            \begin{center}
                \textbf{INSTRUÇÕES PARA A REALIZAÇÃO DESTA AVALIAÇÃO:}
            \end{center}
            \begin{itemize}[leftmargin=*]
                \item Esta é a primeira avaliação da disciplina Sem disciplina, 
                    e engloba o conteúdo referente Sem tag.
                \item Esta avaliação contém 20 questões objetivas, \textbf{todas obrigatórias}.
                \item \textbf{Leia atentamente cada questão!} As questões objetivas terão uma,
                    e apenas uma única, resposta a ser marcada.
                \item \textbf{Desligue o celular} antes de começar. A avaliação é
                    \textbf{individual} e \textbf{sem consulta}.
                \item As questões podem ser respondidas, na prova, com lápis ou caneta. Ao final
                    da prova \textbf{{VOCÊ DEVERÁ PREENCHER O GABARITO COM CANETA
                    DE COR PRETA}} \textbf{\underline{OU AZUL}} para que seja feita a correção
                    automatizada através da leitura do padrão de respostas marcado no
                    gabarito.
                \item \textbf{CUIDADO ao preencher o gabarito} pois eles são \textbf{INDIVIDUAIS
                    e NOMEADOS} (cada estudante tem um gabarito próprio específico, com um
                    código de identificação único). \textbf{Questões rasuradas são anuladas}
                    pelo software de leitura do gabarito.
                \item Ao final da prova, \textbf{DEVOLVA A PROVA E O GABARITO} para o professor.
                \item Siga todas as normas de \textbf{Integridade Acadêmica} da disciplina.
                    Alunos que forem flagrados com qualquer espécie de ``cola'' ou trocando
                    informações com outros alunos terão suas avaliações recolhidas, as notas
                    zeradas, e a situação será encaminhada para a coordenação para a aplicação
                    das penalidades previstas pela universidade.
                \item Com 90 minutos de avaliação você terá tempo de até 4,min por questão,
                    em média.
                \item Boa avaliação!
            \end{itemize}
            \vspace{2cm}
        }}}
    \end{center}
    \newpage

    %%%%%%%%%%%%%%%%%%%%%%%%%%%%%%%%%%%%%%%%%%%%%%%%%%%%%%%%%%%%%%%%%%%%%%%%%%%%%%%%
    %%% Inicia o bloco de questões:
    \begin{questions}

    \newpage
    \question

``Um modelo de dados é uma definição lógica, abstrata e auto-contida

dos \circled{1}, \circled{2} e \circled{3} que, juntos, constituem a \circled{4}

com a qual os usuários interagem. Os \circled{1} nos permitem modelar

a \circled{5}; os \circled{2} nos permitem modelar

seu \circled{6}. Os \circled{3} nos permitem modelar sua \circled{7}.'



As palavras/termos que completam a frase acima, são, conforme os números:


\begin{choices}
\choice 1 = demais conceitos; 2 = objetos; 3 = operadores; 4 = estrutura de dados; 5 = máquina abstrata; 6 = integridade; 7 = comportamento.
\choice 1 = objetos; 2 = operadores; 3 = máquina abstrata; 4 = demais conceitos; 5 = estrutura de dados; 6 = integridade; 7 = comportamento.
\choice 1 = objetos; 2 = máquina abstrata; 3 = demais conceitos; 4 = operadores; 5 = estrutura de dados; 6 = comportamento; 7 = integridade.
\choice 1 = operadores; 2 = objetos; 3 = demais conceitos; 4 = máquina abstrata; 5 = estrutura de dados; 6 = comportamento; 7 = integridade.
\choice 1 = objetos; 2 = operadores; 3 = demais conceitos; 4 = máquina abstrata; 5 = estrutura de dados; 6 = comportamento; 7 = integridade.
\end{choices}

\question

Considere as seguintes afirmativas sobre as linguagens usadas para análise sintática:

\begin{enumerate}[label=\Roman*.]

    \item A classe LL(1) não aceita linguagens com produções que apresentem recursões diretas à esquerda (ex. \(L \rightarrow La\)) mas aceita linguagens com recursões indiretas (ex. \(L \rightarrow Ra, R \rightarrow Lb\)).

    \item A linguagem LR(1) reconhece a mesma classe de linguagens que LALR(1).

    \item A linguagem SLR(1) reconhece uma classe de linguagens maior que LR(0).

\end{enumerate}

Selecione a afirmativa correta:
\begin{choices}
\choice As afirmativas I e III são falsas
\choice As afirmativas II e III são verdadeiras
\choice As afirmativas I e III são verdadeiras
\choice Apenas a afirmativa III é verdadeira
\choice As afirmativas I e II são verdadeiras
\end{choices}

\question

Ao analisar um novo sistema de gerenciamento de bancos de dados que surgiu no

mercado, você percebe que a entrada às operações do modelo de dados são tabelas

mas as saídas nunca são tabelas, são sempre outras estruturas. O vendedor desse

SGBD afirma que ele é um SGBD relacional. Pode-se afirmar que:


\begin{choices}
\choice Nenhuma das alternativas anteriores.
\choice Esse SGBD é relacional pois consegue trabalhar com tabelas para a entrada das operações do modelo de dados. A saída não é importante.
\choice Esse SGBD é relacional pois o vendedor afirmou que é.
\choice Esse SGBD não é relacional pois, para ser verdadeiramente relacional, as operações do modelo de dados não podem trabalhar recebendo tabelas como entrada.
\choice Esse SGBD não é relacional pois no modelo relacional a única estrutura percebida pelos usuários são tabelas.
\end{choices}

\question

Uma árvore binária é declarada em C como



\begin{verbatim}

typedef struct no *apontador;

struct no {

    int valor;

    apontador esq, dir;

};

\end{verbatim}



onde \texttt{esq} e \texttt{dir} representam ligações para os filhos esquerdo e direito de um nó da árvore, respectivamente. Qual das seguintes alternativas é uma implementação correta da operação que inverte as posições dos filhos esquerdo e direito de um nó \texttt{p} da árvore, onde \texttt{t} é um apontador auxiliar.


\begin{choices}
\choice \texttt{p->dir = t;} \\

          \texttt{p->esq = p->dir;} \\

          \texttt{p->dir = t;}
\choice \texttt{t = p;} \\

          \texttt{p->esq = p->dir;} \\

          \texttt{p->dir = p->esq;}
\choice \texttt{p->esq = p->dir;} \\

          \texttt{t = p->esq;} \\

          \texttt{p->dir = t;}
\choice \texttt{t = p->dir;} \\

          \texttt{p->esq = p->dir;} \\

          \texttt{p->dir = t;}
\choice \texttt{t = p->dir;} \\

          \texttt{p->dir = p->esq;} \\

          \texttt{p->esq = t;}
\end{choices}

\question

(ENADE 2014) Considere o seguinte banco de dados ilustrado no diagrama relaconal

abaixo:

\begin{figure}[H]

  \centering

  \caption{Estados e fornecedores}

  \label{fig:estados}

  %\fbox{

    \includegraphics[width=0.5\linewidth]{static/imgs/defaul.png}

  %}\\

  %\footnotesize{Fonte: xxxx.}

\end{figure}

A expressão SQL que obtém os nomes dos estados para os quais não há fornecedores

cadastrados é:


\begin{choices}
\choice \begin{verbatim}SELECT e.nome

FROM estado e,

FROM fornecedor f

WHERE e.uf = f.uf;

\end{verbatim}


\choice \begin{verbatim}SELECT e.uf

FROM estado e

WHERE e.nome NOT IN (SELECT f.uf

                     FROM fornecedor f);

\end{verbatim}
\choice \begin{verbatim}SELECT e.nome

FROM estado e

WHERE e.uf IN (SELECT f.uf

               FROM fornecedor f);

\end{verbatim}


\choice \begin{verbatim}SELECT e.nome

FROM estado e

WHERE e.uf NOT IN (SELECT f.uf

                   FROM fornecedor f);

\end{verbatim}


\choice \begin{verbatim}SELECT e.nome

FROM estado e,

FROM fornecedor f

WHERE e.nome = f.uf;

\end{verbatim}


\end{choices}

\question

Uma integração de Sistemas Computacionais formando uma rede, tipicamente é implementada através da instalação de uma Arquitetura de Rede, que é composta de camadas e protocolos, em cada um dos elementos que compõem esta rede. Considere que estações “conversam” quando aplicações de usuários conseguem comunicar-se, sintática e semanticamente, através da Rede de Computadores. Baseados nesta premissa e em todos os conceitos associados à implementação e utilização das redes de computadores podemos afirmar como certo:


\begin{choices}
\choice Nenhuma delas é uma afirmação correta.
\choice Computadores com arquiteturas de redes diferentes podem “conversar” através de um \textit{gateway} ou conversor de protocolos.
\choice Computadores com arquiteturas diferentes podem “conversar” através de multiplexadores.
\choice Computadores com arquiteturas de rede parecidas conseguem “conversar”.
\choice Computadores com arquiteturas de redes diferentes conseguem “conversar”.
\end{choices}

\question

Em relação ao teste de software, qual das afirmações a seguir é INCORRETA: 
\begin{choices}
\choice Os dados compilados quando a atividade de teste é levada a efeito proporcionam uma boa indicação da confiabilidade do software e alguma indicação da qualidade do software como um todo. 
\choice Um bom caso de teste é aquele que tem uma elevada probabilidade de revelar um erro ainda não descoberto. 
\choice A atividade de teste é o processo de executar um programa com a intenção de demonstrar a ausência de erros. 
\choice O processo de depuração é a parte mais imprevisível do processo de teste pois um erro pode demorar uma hora, um dia ou um mês para ser diagnosticado e corrigido. 
\choice Um teste bem sucedido é aquele que revela um erro ainda não descoberto. 
\end{choices}

\question

A arquitetura ilustrada na figura abaixo é um exemplo típico de qual sistema?

(CU = unidade de controle; IS = \ingles{stream} de instruções; PU = unidade de

processamento; DS = \ingles{stream} de dados; LM = memória local)

\begin{figure}[H]

\begin{center}

\includegraphics[width=0.5\linewidth]{static/imgs/defaul.png}

\end{center}

\end{figure}

\vspace{-1cm}
\begin{choices}
\choice Processadores multicore
\choice Clusters
\choice Multiprocessadores simétricos
\choice Processadores seqüenciais
\choice Acesso não uniforme à memória
\end{choices}

\question

Um compilador detecta:
\begin{choices}
\choice todos os erros citados acima
\choice erros de sintaxe do programa 
\choice erros que podem ocorrer durante a execução do programa
\choice erros aritméticos
\choice erros nos resultados gerados pelo programa
\end{choices}

\question

Seja \( f : \mathbb{R} \to \mathbb{R} \) derivável. Se existem \( a, b \in \mathbb{R} \) tal que \( f(a)f(b) < 0 \) e \( f(x) \neq 0 \) para todo \( x \in (a, b) \), podemos afirmar que no intervalo \( (a, b) \) a equação \( f(x) = 0 \) tem:
\begin{choices}
\choice uma raiz imaginária
\choice somente raízes imaginárias
\choice nenhuma raiz real
\choice uma única raiz real
\choice duas raízes reais
\end{choices}

\question

Avalie o banco de dados de professores e alunos, mostrado a seguir. Esse banco

de dados pretende armazenar todas as turmas, sendo que um professor pode dar

aula para vários alunos ao mesmo tempo, e um mesmo aluno pode ter aula com

diferentes professores ao mesmo tempo.

\begin{figure}[H]

  \centering

  \caption{Professores, turmas e alunos}

  \label{fig:professores}

  %\fbox{

    \includegraphics[width=0.5\linewidth]{static/imgs/defaul.png}

  %}\\

  %\footnotesize{Fonte: xxxx.}

\end{figure}

Esse banco de dados resolveráo problema de armazenar os dados das turmas?
\begin{choices}
\choice Não, pois as cardinalidades e os sentidos dos relacionamentos entre 'professores' e 'turmas', e entre 'turmas' e 'alunos' não cria um relacionamento N:N entre 'professores' e 'alunos'.
\choice Sim, pois o diagrama mostra claramente que um professor pode ter zero ou mais turmas, e cada aluno pode participar de zero ou mais turmas também.
\choice Sim, pois as cardinalidades e os sentidos dos relacionamentos entre 'professores' e 'alunos' indica, no final, um relacionamento N:N entre essas duas tabelas.
\choice Não é possível determinar somente avaliando o diagrama exibido, seriam necessárias mais informações do dicionário de dados.
\choice Não, pois o relacionamento entre 'turmas' e 'professores' está errado: a PK 'cod\_turma' deveria ter ido para a tabela 'professores' como uma FK, mas no diagrama está ao contrário.
\end{choices}

\question

Em um banco de dados relacional, o que é uma chave primária?


\begin{choices}
\choice É qualquer coluna ou conjunto de colunas que podem servir como índice
\choice É denominada por FK
\choice É uma validação que evita o cadastro de dados errados
\choice É um atributo (ou um conjunto de atributos) que permite identificar única e exclusivamente cada registro da tabela
\choice Nenhuma das respostas acima
\end{choices}

\question

O número de \textit{strings binárias} de comprimento 7 e contendo um par de zeros consecutivos é


\begin{choices}
\choice 95
\choice 91
\choice 94
\choice 90
\choice 92
\end{choices}

\question

Considere o seguinte trecho de programa:\\

1. \( i := 1; \)\\

2. while \( i \leq n \) do\\

   begin\\

3. \( sum := sum + a[i]; \)\\

4. \( i := i + 1; \)\\

   end;\\



Considere que:\\

- \( I \) representa a inicialização da variável \( i := 1 \) na linha 1;\\

- \( T \) representa o teste da linha 2;\\

- \( A \) representa os comandos da linha 3;\\

- \( P \) representa o incremento na linha 4.\\



Qual é a expressão regular que representa todas as sequências de passos possíveis de serem executados por este trecho de programa?
\begin{choices}
\choice \( I(TAP)^+ \)
\choice \( I(T(AP)^*)^+ \)
\choice \( I(T(AP)^*)^* \)
\choice \( IT^+A^*P^* \)
\choice \( I(TAP)^* \)
\end{choices}

\question Assinale a alternativa \textbf{incorreta}:
\begin{choices}
\choice Um SGBD não tem interface gráfica
\choice O modelo relacional de dados foi criado por Edgar Frank Cood
\choice O MariaDB é um SGBD
\choice No modelo relacional, tudo o que o usuário percebe são tabelas
\choice Modelo de dados é a mesma coisa que o diagrama relacional
\end{choices}

\question

Considere as seguintes tabelas em uma base de dados relacional:\\

\textbf{Departamento (CodDepto, NomeDepto)}\\

\textbf{Empregado (CodEmp, NomeEmp, CodDepto)}\\

Deseja-se obter uma tabela na qual cada linha é a concatenação de uma linha da tabela Departamento com uma linha da tabela de Empregado. Caso um departamento não possua empregados, sua linha no resultado deve conter vazio (NULL) nos campos referentes ao empregado. A operação de álgebra relacional que deve ser aplicada para combinar estas duas tabelas é:
\begin{choices}
\choice Divisão
\choice Junção externa
\choice União
\choice Junção interna
\choice Projeção
\end{choices}

\question

Dentre as definições a seguir, ligadas ao conceito de visões do modelo relacional, qual delas é \textbf{INCORRETA}?


\begin{choices}
\choice Uma visão relacional é uma relação virtual, derivada de relações base a partir da especificação de operações da álgebra relacional.
\choice Uma visão é útil por representar uma percepção particular do banco de dados, compartilhado por muitos aplicativos.
\choice Programas aplicativos do banco de dados podem ser executados sobre visões de relações da base de dados.
\choice O gerenciamento de visões envolve a conversão da consulta do usuário sobre as visões para a consulta sobre as relações base.
\choice Uma visão relacional é uma relação virtual que nunca é materializada.
\end{choices}

\question

Para que serve a segmentação de um processador (\textit{pipelining})?
\begin{choices}
\choice aumentar a velocidade do relógio
\choice simplificar a implementação do processador
\choice simplificar o conjunto de instruções
\choice reduzir o número de instruções estáticas nos programas
\choice permitir a execução de mais de uma instrução por ciclo de relógio
\end{choices}

\question

Uma definição lógica, abstrata e auto-contida dos objetos que modelam a

estrutura de armazenamento de dados, dos operadores que modelam o comportamento

e dos demais conceitos que modelam a integridade, e que, juntos, constituem a

máquina abstrata com a qual os usuários interagem é chamada de:


\begin{choices}
\choice Banco de dados
\choice Nenhuma das respostas acima
\choice Modelo entidade-relacionamento
\choice Modelo de dados
\choice Modelo relacional
\end{choices}

\question

Na criptografia com chave pública:


\begin{choices}
\choice O sigilo é obtido através da codificação com a chave pública do destinatário e decifragem com a chave privada do destinatário.
\choice Para assinar digitalmente uma mensagem codifica-se a mesma com a chave pública do destinatário e esta é decifrada com a chave privada do destinatário.
\choice Para assinar digitalmente uma mensagem codifica-se a mesma com a chave pública do remetente e esta é decifrada com a chave privada do destinatário.
\choice O sigilo é obtido através da codificação com a chave privada do destinatário e decifragem com a chave pública do destinatário.
\choice O sigilo é obtido através da codificação com a chave privada do remetente e decifragem com a chave pública do destinatário.
\end{choices}



    %%%%%%%%%%%%%%%%%%%%%%%%%%%%%%%%%%%%%%%%%%%%%%%%%%%%%%%%%%%%%%%%%%%%%%%%%%%%%%%%
    %%% Finaliza o bloco de questões:
    \end{questions}
    \end{document}
    